\section{Performance e Otimizacao}

\subsection{Monitoramento de Performance}

\subsubsection{Metricas Essenciais}

\begin{lstlisting}[caption=Queries para monitoramento]
-- CPU e memoria por query
SELECT 
    query,
    calls,
    total_time,
    mean_time,
    stddev_time,
    rows,
    100.0 * shared_blks_hit / nullif(shared_blks_hit + shared_blks_read, 0) AS hit_percent
FROM pg_stat_statements 
ORDER BY mean_time DESC 
LIMIT 10;

-- Bloqueios ativos
SELECT 
    pg_stat_activity.pid,
    pg_stat_activity.usename,
    pg_stat_activity.query,
    pg_stat_activity.state,
    pg_locks.mode,
    pg_locks.locktype,
    pg_locks.relation::regclass
FROM pg_stat_activity
JOIN pg_locks ON pg_stat_activity.pid = pg_locks.pid
WHERE pg_stat_activity.state = 'active';

-- Conexoes ativas por usuario
SELECT 
    usename,
    count(*) as active_connections,
    max(state) as max_state
FROM pg_stat_activity 
WHERE state != 'idle'
GROUP BY usename
ORDER BY active_connections DESC;
\end{lstlisting}

\subsubsection{Dashboard de Monitoramento}

\begin{lstlisting}[caption=View consolidada para dashboard]
CREATE VIEW dashboard_performance AS
WITH estatisticas_gerais AS (
    SELECT 
        (SELECT count(*) FROM usuarios WHERE ativo = true) as usuarios_ativos,
        (SELECT count(*) FROM transacoes WHERE data_transacao = CURRENT_DATE) as transacoes_hoje,
        (SELECT count(*) FROM contas WHERE ativa = true) as contas_ativas,
        (SELECT pg_size_pretty(pg_database_size('finops'))) as tamanho_bd
),
performance_queries AS (
    SELECT 
        count(*) as queries_executadas,
        avg(mean_time) as tempo_medio,
        max(mean_time) as tempo_maximo
    FROM pg_stat_statements
    WHERE calls > 0
),
uso_recursos AS (
    SELECT 
        round(100.0 * sum(shared_blks_hit) / nullif(sum(shared_blks_hit + shared_blks_read), 0), 2) as cache_hit_ratio,
        count(*) as conexoes_ativas
    FROM pg_stat_activity
    WHERE state = 'active'
)
SELECT 
    eg.*,
    pq.queries_executadas,
    pq.tempo_medio,
    pq.tempo_maximo,
    ur.cache_hit_ratio,
    ur.conexoes_ativas,
    CURRENT_TIMESTAMP as ultima_atualizacao
FROM estatisticas_gerais eg, performance_queries pq, uso_recursos ur;
\end{lstlisting}

\subsection{Otimizacoes de Query}

\subsubsection{Analise de Queries Lentas}

\begin{lstlisting}[caption=Identificacao de queries problematicas]
-- Habilitar log de queries lentas
-- No postgresql.conf:
-- log_min_duration_statement = 1000  # log queries > 1s
-- log_statement = 'all'
-- log_duration = on

-- Query para analisar queries mais lentas
SELECT 
    substring(query, 1, 50) as query_inicio,
    calls,
    total_time,
    mean_time,
    rows,
    100.0 * shared_blks_hit / nullif(shared_blks_hit + shared_blks_read, 0) AS hit_percent
FROM pg_stat_statements 
WHERE mean_time > 100  -- queries com media > 100ms
ORDER BY mean_time DESC;

-- Analisar plano de execucao de query especifica
EXPLAIN (ANALYZE, BUFFERS, FORMAT JSON) 
SELECT t.*, c.nome as categoria_nome
FROM transacoes t
JOIN categorias c ON t.categoria_id = c.id
WHERE t.usuario_id = 'uuid_exemplo'
  AND t.data_transacao >= '2024-01-01'
ORDER BY t.data_transacao DESC;
\end{lstlisting}

\subsubsection{Otimizacoes Especificas}

\begin{lstlisting}[caption=Queries otimizadas para casos comuns]
-- Query otimizada para extrato mensal
CREATE INDEX IF NOT EXISTS idx_transacoes_usuario_data_categoria 
    ON transacoes(usuario_id, data_transacao DESC, categoria_id)
    WHERE status = 'confirmada';

-- View materializada para resumos mensais rapidos
CREATE MATERIALIZED VIEW mv_resumo_mensal AS
SELECT 
    usuario_id,
    DATE_TRUNC('month', data_transacao) as mes,
    categoria_id,
    SUM(CASE WHEN tipo = 'receita' THEN valor ELSE 0 END) as receitas,
    SUM(CASE WHEN tipo = 'despesa' THEN valor ELSE 0 END) as despesas,
    COUNT(*) as total_transacoes
FROM transacoes
WHERE status = 'confirmada'
GROUP BY usuario_id, DATE_TRUNC('month', data_transacao), categoria_id;

-- Indice na view materializada
CREATE UNIQUE INDEX idx_mv_resumo_mensal 
    ON mv_resumo_mensal(usuario_id, mes, categoria_id);

-- Query otimizada usando a view
SELECT 
    c.nome as categoria,
    sum(rm.receitas) as total_receitas,
    sum(rm.despesas) as total_despesas
FROM mv_resumo_mensal rm
JOIN categorias c ON rm.categoria_id = c.id
WHERE rm.usuario_id = $1
  AND rm.mes >= DATE_TRUNC('month', CURRENT_DATE - INTERVAL '6 months')
GROUP BY c.nome
ORDER BY sum(rm.despesas) DESC;
\end{lstlisting}

\subsection{Manutencao Preventiva}

\subsubsection{Scripts de Manutencao}

\begin{lstlisting}[caption=Rotina de manutencao semanal]
#!/bin/bash
# manutencao_semanal.sh

LOGFILE="/var/log/finops_manutencao.log"

echo "$(date): Iniciando manutencao semanal" >> $LOGFILE

# 1. VACUUM ANALYZE em tabelas principais
echo "$(date): Executando VACUUM ANALYZE" >> $LOGFILE
psql -d finops -c "VACUUM ANALYZE usuarios;"
psql -d finops -c "VACUUM ANALYZE grupos;"
psql -d finops -c "VACUUM ANALYZE contas;"
psql -d finops -c "VACUUM ANALYZE categorias;"
psql -d finops -c "VACUUM ANALYZE transacoes;"

# 2. Reindex tabelas com muita atividade
echo "$(date): Reindexando tabelas" >> $LOGFILE
psql -d finops -c "REINDEX TABLE transacoes;"
psql -d finops -c "REINDEX TABLE audit.audit_logs;"

# 3. Atualizar estatisticas
echo "$(date): Atualizando estatisticas" >> $LOGFILE
psql -d finops -c "ANALYZE;"

# 4. Refresh das views materializadas
echo "$(date): Atualizando views materializadas" >> $LOGFILE
psql -d finops -c "REFRESH MATERIALIZED VIEW CONCURRENTLY mv_resumo_mensal;"
psql -d finops -c "REFRESH MATERIALIZED VIEW CONCURRENTLY resumo_financeiro_mensal;"

# 5. Limpeza de dados antigos
echo "$(date): Limpando dados antigos" >> $LOGFILE
psql -d finops -c "DELETE FROM audit.audit_logs WHERE timestamp_operacao < CURRENT_DATE - INTERVAL '90 days';"
psql -d finops -c "DELETE FROM sessoes WHERE data_expiracao < CURRENT_TIMESTAMP;"

# 6. Verificar integridade
echo "$(date): Verificando integridade" >> $LOGFILE
psql -d finops -c "SELECT count(*) as usuarios_inconsistentes FROM usuarios WHERE email IS NULL OR email = '';"

echo "$(date): Manutencao semanal concluida" >> $LOGFILE
\end{lstlisting}

\subsubsection{Monitoramento Automatico}

\begin{lstlisting}[caption=Script de alerta de performance]
#!/bin/bash
# alertas_performance.sh

DB_NAME="finops"
EMAIL_ADMIN="admin@finops.com"
LIMITE_CONEXOES=50
LIMITE_TEMPO_QUERY=5000  # 5 segundos
LIMITE_CACHE_HIT=95      # 95%

# Verificar numero de conexoes
CONEXOES_ATIVAS=$(psql -d $DB_NAME -t -c "SELECT count(*) FROM pg_stat_activity WHERE state = 'active';")

if [ $CONEXOES_ATIVAS -gt $LIMITE_CONEXOES ]; then
    echo "ALERTA: Muitas conexoes ativas ($CONEXOES_ATIVAS)" | \
    mail -s "FinOps Performance Alert" $EMAIL_ADMIN
fi

# Verificar queries lentas
QUERIES_LENTAS=$(psql -d $DB_NAME -t -c "
    SELECT count(*) FROM pg_stat_statements 
    WHERE mean_time > $LIMITE_TEMPO_QUERY AND calls > 10;
")

if [ $QUERIES_LENTAS -gt 0 ]; then
    echo "ALERTA: $QUERIES_LENTAS queries com performance ruim" | \
    mail -s "FinOps Query Performance Alert" $EMAIL_ADMIN
fi

# Verificar cache hit ratio
CACHE_HIT=$(psql -d $DB_NAME -t -c "
    SELECT round(100.0 * sum(blks_hit) / nullif(sum(blks_hit + blks_read), 0), 2)
    FROM pg_stat_database WHERE datname = '$DB_NAME';
")

if (( $(echo "$CACHE_HIT < $LIMITE_CACHE_HIT" | bc -l) )); then
    echo "ALERTA: Cache hit ratio baixo ($CACHE_HIT%)" | \
    mail -s "FinOps Cache Performance Alert" $EMAIL_ADMIN
fi

# Verificar espaco em disco
ESPACO_BD=$(psql -d $DB_NAME -t -c "SELECT pg_database_size('$DB_NAME');")
ESPACO_GB=$((ESPACO_BD / 1024 / 1024 / 1024))

if [ $ESPACO_GB -gt 10 ]; then
    echo "AVISO: Banco de dados com ${ESPACO_GB}GB" | \
    mail -s "FinOps Storage Warning" $EMAIL_ADMIN
fi
\end{lstlisting}

\subsection{Configuracoes de Performance}

\subsubsection{Parametros Otimizados}

\begin{lstlisting}[caption=Configuracoes recomendadas do PostgreSQL]
# postgresql.conf - Configuracoes para servidor com 4GB RAM

# MEMORIA
shared_buffers = 1GB                    # 25% da RAM
effective_cache_size = 3GB              # 75% da RAM
work_mem = 16MB                         # Para operacoes complexas
maintenance_work_mem = 256MB            # Para VACUUM, CREATE INDEX

# CHECKPOINT E WAL
wal_buffers = 16MB
checkpoint_completion_target = 0.9
max_wal_size = 2GB
min_wal_size = 512MB
checkpoint_timeout = 10min

# PLANNER
effective_io_concurrency = 200         # Para SSD
random_page_cost = 1.1                 # Para SSD
default_statistics_target = 100

# CONEXOES
max_connections = 100
shared_preload_libraries = 'pg_stat_statements'

# LOG E MONITORAMENTO
log_min_duration_statement = 1000      # Log queries > 1s
log_checkpoints = on
log_connections = on
log_disconnections = on
log_lock_waits = on
log_statement = 'mod'
log_line_prefix = '%t [%p]: [%l-1] user=%u,db=%d,app=%a,client=%h '

# AUTO VACUUM
autovacuum = on
autovacuum_max_workers = 3
autovacuum_naptime = 20s
\end{lstlisting}
